\section{Types of SAT Encodings}

There are three different types of SAT Encodings planned in conjure oxide. Of these, only Logarithmic Encodings have been implemented thus far. The three types are these:

\begin{itemize}
    \item Direct Encodings
    \item Logarithmic Encodings
    \item Order Encodings
\end{itemize}

Each type of encoding has pros and cons, and a different one may be selected based on the type of constraint problem.

\subsection{Logarithmic Encoding}

The base principle is quite simple: encode an integer as a bitvector. This allows us to represent integers as a series of boolean constraints -- one for each bit. 

For example, the integer `6' can be represented in binary as '0110'. We can then represent this as ($P = 0\lor1, Q = {0\lor1}, R = {0\lor1}, S = {0\lor1}$). The connection that is missing, however, is that this isn't actually a representation of a constraint problem, but of a solution to a problem.  

\subsection{Direct Encoding}

Direct encodings are the most straightforward type of encoding -- it involves creating a vector of boolean variables, one corresponding to each member of the domain. Only one of these variables can be true at a time, and it is the one corresponding to the value of the integer. 

\subsection{Order Encoding}

Order follows the same principle as direct encoding, but instead of each boolean variable `specifying' a value in the way that direct encodings do, each bit specifies whether the integer corresponding to it is less than or equal to the integer variable's value.

\subsection{Why want to add multiple types of encoding?}

Only logarithmic encodings are currently implemented in conjure-oxide. We're planning to include other encoding types such as direct and order encodings. This is motivated by their potential advantages over the log encoding in some cases.

Direct encodings should perform well for equality-heavy constraints but may scale poorly with larger domains or inequalities. Logarithmic encodings are expected to handle inequalities more efficiently. Order encodings are often viewed as a compromise, potentially balancing these trade-offs.

\subsection{Current implementation status}

As of now (17/12/2025), only Logarithmic Encodings have been implemented in conjure oxide. Future work includes implementing Direct and Order Encodings to provide more options for different types of constraint problems. We also are looking into hybrid approaches that can combine multiple encoding types within a single problem instance to leverage the strengths of each where appropriate.

\subsubsection{Examples illustrating encoding choices}
THIS SUBSECTION IS DRAFT AND NOT GONNA BE A PART OF CURRENT BOOK. IT WOULD BE INCLUDED ONCE WE HAVE ORDER AND DIRECT ENCODINGS RUNNING.
\paragraph{Direct encoding (equality / all-different).}
\begin{verbatim}
find x,y,z : int(1..3)
such that
    x != y,
    y != z,
    x != z
\end{verbatim}

\noindent
\textbf{Why direct works well:}
In a direct encoding, each possible value of a variable is represented explicitly by a Boolean. 
(Here I realised that we don't really have a good explanation of direct encoding earlier. TODO: add one)
The constraint $x \neq y$ can therefore be expressed by a small set of simple clauses that forbid
$x$ and $y$ from taking the same value.
For this example, all constraints are pairwise equalities, so they translate directly into short, local clauses that immediately rule out conflicting assignments.
In contrast, a logarithmic encoding would require reasoning over multiple bits to detect equality,
resulting in weaker propagation and more complex clauses.

\paragraph{Logarithmic encoding (arithmetic / inequalities).}
\begin{verbatim}
find x,y : int(0..15)
such that
    x + y <= 12,
    x < y
\end{verbatim}

\noindent
\textbf{Why logarithmic works well:}
Here the constraints are dominated by arithmetic and ordering.
With a logarithmic encoding, $x$ and $y$ are represented using only four bits each, and both
$x < y$ and $x + y \le 12$ can be compiled into standard comparator and adder circuits (which is what I think is happening in our conjure-oxide implementation. I need Henry hlep here).
This results in a relatively small number of Boolean variables and avoids enumerating forbidden
value pairs.
A direct encoding would require clauses for many combinations of $(x,y)$ whose sum exceeds $12$,
leading to a much larger and less structured SAT instance.

\paragraph{Order encoding (bounds / ranges / cardinality).}
\begin{verbatim}
find x,y,z : int(1..10)
such that
    x <= 4,
    y <= 4,
    z <= 4
\end{verbatim}

\noindent
\textbf{Why order works well:}
Order encodings are particularly effective for bound constraints.
Each constraint of the form $x \le k$ can be enforced by setting a single Boolean variable,
which immediately propagates through the ordering structure of the encoding.
As a result, the SAT solver can prune large parts of the search space early.
In comparison, both direct and logarithmic encodings would require multiple clauses to express
the same bound, making propagation less immediate.

\paragraph{Mixed illustration.}
\begin{verbatim}
find x,y : int(1..8)
such that
    x != y,
    x + y <= 9
\end{verbatim}

\noindent
\textbf{Why this is challenging:}
This example combines equality and arithmetic constraints.
The inequality $x \neq y$ is naturally handled by a direct encoding, while the arithmetic constraint
$x + y \le 9$ is more naturally handled by a logarithmic encoding.
Any single encoding choice therefore favours one constraint at the expense of the other.
This motivates the need for supporting multiple encodings, and potentially hybrid approaches,
where different variables or constraints are encoded using different strategies.